\section{Proof of the triangular spectral limit}\label{triapp:all}
This appendix proves Theorem~\ref{tri:main}, retaining the zero-counting
and centering estimates that determine its normalization.
Throughout, $c=e^{-\gamma}$.
\subsection{An entire function with a measurable supply of zeros}
Put
\[
 \mathcal E(z)=\sum_{k\ge0}K(k)\frac{z^k}{k!},\qquad
 \kappa(t)=-\log K(t)\quad(t\ge0).
\]
The more accurate prime estimates needed below are
\begin{equation}\label{tri:kappa}
 \begin{split}
 \kappa'(t)&=\log\log t+\gamma+O((\log t)^{-2}),\\
 \kappa''(t)&=\frac{1+O((\log t)^{-1})}{t\log t},\\
 |\kappa'''(t)|&\ll(t^2\log t)^{-1}.
 \end{split}
\end{equation}
Here and below derivative asymptotics concern $t\to\infty$.

We explain the extra precision in the first line, since an error of
order $1/\log t$ would change the final location.  Differentiation
gives
\[
 \kappa'(t)=\sum_p\left\{-\log(1-1/p)-\frac1{p+t-1}\right\}.
\]
After splitting at $p=t$, Mertens' product theorem supplies
$\log\log t+\gamma$ with error $o((\log t)^{-2})$.
Replacing $p+t-1$ by $p+t$, and the logarithm by $1/p$ above $t$,
costs $O(1/(t\log t))$.  The remaining prime sum is
\[
 -\sum_{p\le t}\frac1{p+t}
       +\sum_{p>t}\frac{t}{p(p+t)}.
\]
Partial summation with the classical prime-number-theorem error
turns its part on $t^{1/2}\le p\le t^{3/2}$ into the integral
of $h(v)/(\log t+\log v)$, where
\[
 h(v)=
 \begin{cases}
 -(1+v)^{-1},&0<v\le1,\\
 [v(1+v)]^{-1},&v>1.
 \end{cases}
\]
The omitted ranges are $O(t^{-1/2}/\log t)$.
Since $\int_0^\infty h(v)\,dv=0$ and
$\int_0^\infty |h(v)\log v|\,dv<\infty$, the correction is
$O((\log t)^{-2})$.  This proves the first line of
\eqref{tri:kappa}.  Applying the same partial summation to
$\sum_p(p+t-1)^{-2}$ proves the second; differentiating once more
and bounding $\sum_p(p+t-1)^{-3}$ proves the third.

In particular $K(k)^{1/k}\to0$.  For every $\epsilon>0$,
\[
 |\mathcal E(z)|\le C_\epsilon e^{\epsilon|z|}.
\]
Also $J_N(z/N)\to\mathcal E(z)$ locally uniformly, since its $k$th
coefficient tends to $K(k)/k!$ and is bounded by $1/k!$.
The negative real roots of $J_N$ therefore give
\begin{equation}\label{tri:product}
 \mathcal E(z)=\prod_{j\ge1}(1+z/\rho_j),\qquad
 \rho_j>0,\qquad \sum_j\rho_j^{-1}=1.
\end{equation}
Here is a direct way to check the representation and its missing
exponential factor.  Write
$J_N(z/N)=\prod_j(1+\alpha_{N,j}z)$ with
$\alpha_{N,j}\ge0$ and $\sum_j\alpha_{N,j}=1$.
Sort these numbers and take coordinate limits along a subsequence.
The small factors, using
$\log(1+\alpha z)=\alpha z+O_z(\alpha^2)$, contribute an exponential
$e^{az}$ with $a\ge0$; the remaining factors give the product in
\eqref{tri:product}.  A positive $a$ contradicts the zero exponential
type bound on the positive axis.  Hence $a=0$, and
$\mathcal E'(0)=K(1)=1$ gives the reciprocal sum.
There must be infinitely many factors, since every coefficient of
$\mathcal E$ is positive.  In particular $\mathcal E(-z)$ belongs
to the Laguerre--Pólya class and has only positive zeros.

\begin{lemma}[Arithmetic zero density]\label{tri:zeros}
With $n(r)=\#\{j:\rho_j\le r\}$,
\[
 n(r)\sim c\,\frac r{\log^2r},\qquad
 \rho_n\sim c^{-1}n\log^2n.
\]
\end{lemma}
\begin{proof}
For $r>0$, define a random nonnegative integer $L_r$ by
\[
 \Prob(L_r=k)=\frac{K(k)r^k/k!}{\mathcal E(r)}.
\]
Choose $m>0$ from $r=me^{\kappa'(m)}$; the right side is strictly
increasing from zero to infinity.  If $P$ is Poisson with mean $m$,
put $\Delta=P-m$ and
\[
 \beta(k)=\kappa(k)-\kappa(m)-\kappa'(m)(k-m)\ge0.
\]
Convexity of $\kappa$ gives the inequality, and direct cancellation
gives the exact change of measure
\begin{equation}\label{tri:poissontilt}
 \Prob(L_r=k)=
 \frac{e^{-\beta(k)}}{\E e^{-\beta(P)}}\Prob(P=k).
\end{equation}

Write $A=\kappa''(m)$.  For $|\Delta|\le m/2$,
\[
 \beta(P)=\tfrac12A\Delta^2+
 O\left(\frac{|\Delta|^3}{m^2\log m}\right).
\]
The complementary Poisson tail is exponentially small; multiplication
by any of the fixed powers below preserves that bound.
For example, $\kappa(k)=O(k\log(k+2))$ controls $\beta(k)$ there.
Using the Poisson central moments
\[
 \E\Delta^2=m,\quad \E\Delta^3=m,\quad
 \E\Delta^4=3m^2+m,\quad
 \E\Delta^6=15m^3+25m^2+m
\]
and Cauchy--Schwarz for intermediate absolute moments yields
\[
 \begin{aligned}
 \E\beta(P)&=\tfrac12Am+O(m^{-1/2}/\log m),&
 \E\beta(P)^2&=O((\log m)^{-2}),\\
 \E[\Delta^2\beta(P)]
   &=\tfrac12A(3m^2+m)+O(\sqrt m/\log m),&
 \E[\Delta^2\beta(P)^2]&=O(m/(\log m)^2).
 \end{aligned}
\]
Insert $e^{-\beta}=1-\beta+O(\beta^2)$ in
\eqref{tri:poissontilt}.  The mean displacement is
$O(\sqrt m/\log m)$, and the second moment about $m$ is
\[
 \E(L_r-m)^2
 =m-\tfrac12A(2m^2+m)+o(m/\log m)
 =m-\frac m{\log m}+o(m/\log m).
\]
The square of the mean displacement is $O(m/(\log m)^2)$.
Consequently
\begin{equation}\label{tri:vardeficit}
 \Var(L_r)-\E L_r=-\frac m{\log m}(1+o(1)).
\end{equation}

For $g(r)=\log\mathcal E(r)$, differentiation of the probability
weights and of \eqref{tri:product} gives
\[
 g'(r)=\frac{\E L_r}{r},\qquad
 g''(r)=\frac{\Var(L_r)-\E L_r}{r^2}
      =-\sum_j(r+\rho_j)^{-2}.
\]
Thus
\begin{equation}\label{tri:stieltjes}
 \sum_j(r+\rho_j)^{-2}\sim\frac c{r\log^2r}.
\end{equation}
We will also need the more accurate first-derivative estimate
\begin{equation}\label{tri:firstderivative}
 g'(r)=\frac c{\log m}+O((\log m)^{-3}),
 \qquad r=me^{\kappa'(m)},
\end{equation}
which follows from the mean bound and \eqref{tri:kappa}.

For clarity, we give the short Tauberian step that extracts $n(r)$.
Form the positive measure
\[
 d\eta_r(u)=\frac{\log^2r}{cr}\frac{dn(ru)}{(1+u)^2}
 \quad (u\ge0)
\]
and push it forward under $v=(1+u)^{-1}$ to a measure $\lambda_r$
on $[0,1]$.  Its mass tends to one by \eqref{tri:stieltjes}, and
for every $t>0$,
\[
 \int_0^1\frac{d\lambda_r(v)}{[1+(t-1)v]^2}
 =\frac{r\log^2r}{c}\sum_j(tr+\rho_j)^{-2}
 \longrightarrow\frac1t.
\]
The uniform probability measure on $[0,1]$ has these same integrals.
Expanding near $t=1$ determines all moments, which determine a measure
on this compact interval.  Every subsequential limit is therefore
uniform.  It follows that
\[
 \frac{\log^2r}{cr}n(r)
 =\int_{1/2}^1v^{-2}\,d\lambda_r(v)\longrightarrow1.
\]
The discontinuity at $1/2$ has zero limiting mass.
Asymptotic inversion now gives the formula for $\rho_n$.
\end{proof}

\subsection{The transfer theorem and the location constant}
We invoke Campbell and Jalowy
\cite[Corollary 2.3(2)]{finCampbellJalowy}.
For a Laguerre--Pólya function with zero-counting law regularly varying
with index $\alpha\in(0,2)$, their theorem gives a free-stable limit for
suitably scaled and centered Appell roots.  The one-sided index-one
case is included.

Apply it to $F(z)=\mathcal E(-z)$, whose compensated product is
\[
 F(z)=e^{-z}\prod_j(1-z/\rho_j)e^{z/\rho_j}.
\]
Its linear parameter is $1$, its Gaussian parameter is zero,
and all its zeros are positive.  Put
\[
 B_n=\rho_n,\qquad
 S(B)=\sum_j\frac{\rho_j}{\rho_j^2+B^2}.
\]
In the convention of the cited theorem, take limiting location zero.
Its prescribed exponential shift is
\[
 a_n=-\frac{B_n}{n}
   +\frac1n\sum_j
       \frac{(B_n/\rho_j)^3}{1+(B_n/\rho_j)^2}
   =-\frac{B_n}{n}S(B_n),
\]
where the equality uses $\sum_j1/\rho_j=1$.
The Appell polynomial of $F$ is
\[
 A_{n,F}(x)=\sum_{k=0}^n(-1)^kK(k)\binom nkx^{n-k}
          =x^nJ_n(-1/x).
\]
Its roots are $R_{n,j}^{-1}$.  Replacing $F(z)$ by
$e^{-a_nz}F(B_nz/n)$ multiplies these roots by $B_n/n$
and adds $a_n$.  The transfer theorem therefore gives
\begin{equation}\label{tri:appelllimit}
 \frac1n\sum_j\delta_{X_{n,j}}\Longrightarrow\mu_+,\qquad
 X_{n,j}=\frac{B_n}{n}\{R_{n,j}^{-1}-S(B_n)\}.
\end{equation}
The limit has weighted Lévy measure
$\mathbf{1}_{x>0}(1+x^2)^{-1}dx$ and zero location.

To identify its normalization, the free Lévy--Khintchine formula gives
\[
 R_{\mu_+}(u)=\int_0^\infty
                  \frac{u+x}{(1-xu)(1+x^2)}\,dx=-\log(-u).
\]
For instance, at $u=-s<0$ the integrand equals
$x/(1+x^2)-s/(1+sx)$, whose integral is $-\log s$.
Analytic continuation fixes the branch on the lower half-plane.

The remaining step is arithmetic: replace $S(B_n)$ by an explicit
elementary center.  We claim
\begin{equation}\label{tri:centercancellation}
 S(B)-g'(B)=o((\log B)^{-2}).
\end{equation}
Lemma~\ref{tri:zeros} makes the rescaled measure
$(\log^2B/B)\,dn(Bu)$ converge vaguely to $c\,du$, so
\[
 \begin{split}
 &\log^2B\{S(B)-g'(B)\}\\
 &\quad=\int_0^\infty
 \left\{\frac u{1+u^2}-\frac1{1+u}\right\}
               \frac{\log^2B}{B}\,dn(Bu)
 \longrightarrow0.
 \end{split}
\]
Indeed the kernel is bounded at zero, is $O(u^{-2})$ at infinity,
and has integral zero: its antiderivative is
$\tfrac12\log(1+u^2)-\log(1+u)$.
To justify the tails, below $u=B^{-1/2}$ the rescaled mass is
$O(B^{-1/2})$; from there to a fixed small $\epsilon$,
$n(Bu)\ll Bu/\log^2B$ gives $O(\epsilon)$.
Above a fixed $A>1$, the same bound and partial summation give
$O(A^{-1})$.  This proves \eqref{tri:centercancellation}.

Choose $m_n$ with $B_n=m_ne^{\kappa'(m_n)}$.
The zero density gives $m_n\sim n\log n$, and hence
$\log m_n=\log n+\log\log n+o(1)$.
Combining \eqref{tri:firstderivative} and
\eqref{tri:centercancellation} yields
\begin{equation}\label{tri:center}
 \frac c{S(B_n)}=\log n+\log\log n+o(1),\qquad
 \frac{cn}{B_nS(B_n)^2}\to1,\qquad
 \frac n{B_nS(B_n)}\to0.
\end{equation}
The exact identity
\[
 cR_{n,j}-\frac c{S(B_n)}
 =-\frac{cn}{B_nS(B_n)^2}\,
       \frac{X_{n,j}}{1+\frac n{B_nS(B_n)}X_{n,j}}
\]
now completes the transfer.  The empirical laws of $X_{n,j}$ are
tight by \eqref{tri:appelllimit}; on each fixed compact interval
the right side converges uniformly to $-X_{n,j}$.
Thus \eqref{tri:centered} has the reflected limit law of $T$,
with $R$-transform $\log u$.

The limiting density has the following explicit parametrization:
\begin{equation}\label{tri:density}
 x(\theta)=\log\frac{\sin\theta}{\theta}+\theta\cot\theta,\qquad
 f_T(x(\theta))=\frac{\sin^2\theta}{\pi\theta},
 \quad 0<\theta<\pi.
\end{equation}
For completeness, it can be read from the Cauchy-transform
equation $z=1/G(z)+\log G(z)$.
On the support write $G(x+i0)=r e^{-i\theta}$.
The imaginary part forces $r=\sin\theta/\theta$; the real part and
Stieltjes inversion give exactly \eqref{tri:density}.
The function $x(\theta)$ decreases from $1$ to $-\infty$.
Arizmendi and Hasebe \cite[Proposition 4.4]{triArizmendiHasebe}
identify the exponential of this precisely normalized law with
$\operatorname{DH}_1$.
Alternatively, applying the exponential change of variables to
\eqref{tri:density} gives the Dykema--Haagerup density directly.
The continuous mapping theorem proves \eqref{tri:exp}.


